\section{A finite integer-selected real generator with every original jet}\label{sec:constructive-generator}

\subsection{Original objects and the integer selector}

Let $Z=\{\rho_1,\ldots,\rho_r\}$ be a nonempty finite set of actual
zeros of $g=2\xi$, invariant under $\iota\rho=1-\bar\rho$.
Retain each full multiplicity $m_\rho\geq1$, so that
$m_{\iota\rho}=m_\rho$, and put $d=\sum_{\rho\in Z}m_\rho$.
Keep the original local coordinates and reflection:
\begin{equation}
 A_Z=\prod_{\rho\in Z}\mathbb{C}[t_\rho]/(t_\rho^{m_\rho}),\qquad
 t_\rho=s-\rho,\qquad
 (f^\dagger)_{\rho,j}=(-1)^j\overline{f_{\iota\rho,j}},\qquad
 A_{Z,\mathbb{R}}=A_Z^{\dagger}.
 \tag{CG1}\label{cg:original}
\end{equation}
The relation $f=f^\dagger$ is therefore
$f_{\iota\rho,j}=(-1)^j\overline{f_{\rho,j}}$.
No original coordinate or multiplicity is changed below.
For $p=2,3$ write
\begin{equation}
 \lambda_p(s)=p^{s-1/2},\qquad
 a_p(s)=\frac{\lambda_p(s)+\lambda_p(s)^{-1}}2,\qquad
 b_p(s)=\frac{\lambda_p(s)-\lambda_p(s)^{-1}}{2i}.
 \tag{CG2}\label{cg:prime}
\end{equation}
Their images in $A_Z$ are their full Taylor expansions through
$t_\rho^{m_\rho-1}$. The original prime action is
$\Pi_p=M_{p^s}=p^{1/2}M_{\lambda_p}$, not $M_{\lambda_p}$.
The identities $\lambda_p^\dagger=\lambda_p^{-1}$,
$a_p^\dagger=a_p$, $b_p^\dagger=b_p$, and $a_p^2+b_p^2=1$
hold with these exact expressions.

For an integer $c$ set
\begin{equation}
 H_c=a_2+c b_2+c^2a_3+c^3b_3,\qquad
 z_\rho=H_c(\rho),\qquad d_\rho=H_c'(\rho).
 \tag{CG3}\label{cg:generator}
\end{equation}
Here $H_c$ also denotes the entire function in the original variable
$s$; its class is taken only when it is used in $A_Z$.

\begin{theorem}[A bounded finite integer choice]\label{cg:selector-thm}
Put $r_2=\#\{\rho\in Z:m_\rho>1\}$ and
\begin{equation}
 B=3\left(\binom r2+r_2\right).
 \tag{CG4}\label{cg:bound}
\end{equation}
At least one $c\in\{0,1,\ldots,B\}$ has distinct values $z_\rho$
at all distinct roots and $d_\rho\ne0$ at every root of multiplicity greater than one. Every such
$c$ gives $\mathbb{R}[H_c]=A_{Z,\mathbb{R}}$ and $\mathbb{C}[H_c]=A_Z$.
No critical-line or simple-zero hypothesis is needed.
\end{theorem}

\begin{proof}
If both $\lambda_2(\rho)=\lambda_2(\sigma)$ and
$\lambda_3(\rho)=\lambda_3(\sigma)$, then
$(\rho-\sigma)\log2=2\pi i n$ and
$(\rho-\sigma)\log3=2\pi i m$ for integers $n,m$.
When $\rho\ne\sigma$, both are nonzero and elimination gives
$m\log2=n\log3$. Exponentiation gives $2^m=3^n$;
clearing negative exponents contradicts unique prime factorization.
Thus the two original prime values separate distinct roots.
Since $\lambda_p=a_p+ib_p$, the four values
$a_2,b_2,a_3,b_3$ also separate distinct roots.
Consequently, for each unordered pair, the polynomial
\[
 \Delta_{\rho\sigma}(C)=
 (a_2(\rho)-a_2(\sigma))
 +C(b_2(\rho)-b_2(\sigma))
 +C^2(a_3(\rho)-a_3(\sigma))
 +C^3(b_3(\rho)-b_3(\sigma))
\]
is a nonzero complex polynomial of degree at most three.
Also
\[
 a_p'=i(\log p)b_p,\qquad b_p'=-i(\log p)a_p.
\]
Since $a_2^2+b_2^2=1$, the first two coefficients of
\[
 D_\rho(C)=a_2'(\rho)+C b_2'(\rho)
                  +C^2a_3'(\rho)+C^3b_3'(\rho)
\]
cannot both vanish. Thus $D_\rho$ is nonzero of degree at most
three. A nonzero polynomial of degree $q$ has at most $q$ roots:
successive division by $C-c$ at distinct roots proves this by
induction on $q$. Exclude roots of $D_\rho$ only when $m_\rho>1$.
The total number of excluded complex numbers,
and hence of excluded integers, is at most $B$.
The $B+1$ displayed integers therefore contain a permitted one.
The generation assertions follow from the explicit inverse proved
in Theorem~\ref{cg:inverse-thm} below.
\end{proof}

For a completely specified finite selector one may define
\begin{equation}
 \Xi(C)=\prod_{\substack{\rho\in Z\\m_\rho>1}}D_\rho(C)
                 \prod_{\rho<\sigma}\Delta_{\rho\sigma}(C),\qquad
 c=\min\operatorname*{arg\,max}_{0\le n\le B}|\Xi(n)|.
 \tag{CG5}\label{cg:selector}
\end{equation}
The maximum is positive by the proof. This is an exact finite
formula in the supplied zeros, not a replacement by rounded
ordinates or a claim to have certified unspecified numerical inputs.
Derivatives at simple factors are unused. If a nonzero derivative
at every root is additionally wanted, test instead
$0,\ldots,3\binom r2+3r$ and include every $D_\rho$ in $\Xi$.
The sharper bound CG4 suffices for every formula below. An empty
product in CG5 is one.

Every coefficient needed in the construction has the following
explicit expression in the original prime jets. Put
$\ell_p=\log p$ and $x_{p,\rho}=p^{\rho-1/2}$. Then
\begin{equation}
\begin{split}
 h_{\rho,j}:=[t_\rho^j]H_c(\rho+t_\rho)
 =\frac1{2j!}\bigl\{
 &\ell_2^j[(1-ic)x_{2,\rho}
               +(-1)^j(1+ic)x_{2,\rho}^{-1}]\\
 +&\ell_3^j[(c^2-ic^3)x_{3,\rho}
               +(-1)^j(c^2+ic^3)x_{3,\rho}^{-1}]
 \bigr\}.
\end{split}
\tag{CG6}\label{cg:coefficients}
\end{equation}
In particular $h_{\rho,0}=z_\rho$ and $h_{\rho,1}=d_\rho$.
The full reflection identity is
$h_{\iota\rho,j}=(-1)^j\overline{h_{\rho,j}}$.
At a fixed root $z_\rho$ is real and $d_\rho$ is purely imaginary;
at a nonfixed pair the two distinct nodes are conjugate and nonreal.

\subsection{The exact confluent evaluation determinant}

Order the roots as $\rho_1,\ldots,\rho_r$ and order the rows first
by root, then by $j=0,\ldots,m_\rho-1$. Columns are indexed by
$k=0,\ldots,d-1$. Define
\begin{equation}
 E_{(\rho,j),k}=[t_\rho^j]H_c(\rho+t_\rho)^k
       =\frac1{j!}\left.\frac{d^j}{ds^j}H_c(s)^k\right|_{s=\rho}.
 \tag{CG7}\label{cg:evaluation}
\end{equation}
This is the coefficient matrix of
$P\mapsto(P(H_c))_\rho$ in the original $t_\rho$ coordinates.

\begin{theorem}[The determinant, with its orientation and constants]
\label{cg:det-thm}
For the stated row and column orders,
\begin{equation}
 \det E=
 \prod_{\rho\in Z}d_\rho^{m_\rho(m_\rho-1)/2}
 \prod_{1\le a<b\le r}
          (z_{\rho_b}-z_{\rho_a})^{m_{\rho_a}m_{\rho_b}}.
 \tag{CG8}\label{cg:det}
\end{equation}
If row $(\rho,j)$ instead contains the unscaled derivative
$\left.\frac{d^j}{ds^j}H_c(s)^k\right|_{s=\rho}$, its determinant
is CG8 multiplied by
$\prod_\rho\prod_{j=0}^{m_\rho-1}j!$.
\end{theorem}

\begin{proof}
Put $k_\rho(t)=H_c(\rho+t)-z_\rho$ and introduce the matrices
\[
 V_{(\rho,l),k}=\binom kl z_\rho^{k-l},\qquad
 (B_\rho)_{j,l}=[t^j]k_\rho(t)^l
 \quad(0\le j,l<m_\rho).
\]
The first formula is understood as zero for $l>k$; for $l=k$
its value is one, including $z_\rho=0$. Expanding a polynomial
at $z_\rho$ and substituting $k_\rho(t)$ proves
$E=\operatorname{diag}(B_\rho)V$.
Each $B_\rho$ is lower triangular, with diagonal
$1,d_\rho,d_\rho^2,\ldots,d_\rho^{m_\rho-1}$, so its determinant
is the corresponding factor in CG8. When $m_\rho=1$ the block
is $(1)$ and its determinant factor is one, even if $d_\rho=0$.

To compute $\det V$ without omitting its sign, replace the monomial
column basis by the polynomials
\[
 N_{a,j}(X)=\prod_{b<a}(X-z_{\rho_b})^{m_{\rho_b}}
                       (X-z_{\rho_a})^j,
 \qquad 0\le j<m_{\rho_a},
\]
in that order. Their degrees are exactly $0,\ldots,d-1$, and they
are monic. The column change has determinant one. At every earlier
root the whole block of jets vanishes, so the evaluation matrix
is block lower triangular. In its block at $\rho_a$, the diagonal
entry for every $j$ is
$\prod_{b<a}(z_{\rho_a}-z_{\rho_b})^{m_{\rho_b}}$.
That block is itself lower triangular, since $(X-z_{\rho_a})^j$
has no Taylor coefficient of order less than $j$.
Multiplying its diagonal entries over all $a,j$ yields the second
product of CG8 with exactly the displayed orientation.
Finally changing a Taylor-coefficient row into a derivative row
multiplies it by $j!$, which proves the last assertion.
\end{proof}

There is also a fully specified real determinant if required.
Order all fixed roots first, and then consecutive reflection pairs
$(\rho,\iota\rho)$; within each complex factor use increasing $j$.
For a fixed root take real output coordinates $i^j f_{\rho,j}$,
the coefficients in $y_\rho=t_\rho/i$. For a pair take first all
$\Re f_{\rho,j}$ and then all $\Im f_{\rho,j}$, where the second
factor is determined by reflection. The real matrix $E_\mathbb{R}$ on
real polynomial coefficients then satisfies
\begin{equation}
 \det E_\mathbb{R}=
 \left(\prod_{\rho=\iota\rho}i^{m_\rho(m_\rho-1)/2}\right)
 \left(\prod_{\{\rho,\iota\rho\},\,\rho\ne\iota\rho}
       (i/2)^{m_\rho}(-1)^{m_\rho(m_\rho-1)/2}\right)\det E.
 \tag{CG9}\label{cg:real-det}
\end{equation}
This expression is real and nonzero. Indeed the fixed-factor row
change has diagonal $i^j$. At a pair the row-change matrix is
$\left(\begin{smallmatrix}I/2&D/2\\ I/(2i)&-D/(2i)\end{smallmatrix}\right)$,
where $D=\operatorname{diag}((-1)^j)$; its determinant is the
displayed pair factor. On columns $H_c^k$ the rows are real by
CG1, establishing the reality assertion as well as the formula.

\subsection{An inverse that recovers each original Taylor coefficient}

For $m=m_\rho$, let
\begin{equation}
 q_\rho(z)=\sum_{n=1}^{m-1}v_{\rho,n}z^n,
 \qquad k_\rho(q_\rho(z))\equiv z\pmod{z^m}.
 \tag{CG10}\label{cg:inverse-coordinate}
\end{equation}
For $m=1$ this is the zero polynomial in the zero-dimensional
positive-order ideal. For $m\ge2$, all its coefficients are fixed
by the finite recursion
\begin{equation}
 \begin{split}
 v_{\rho,1}&=d_\rho^{-1},\\
 v_{\rho,n}&=-d_\rho^{-1}\sum_{l=2}^{n}h_{\rho,l}
   [z^n]\left(\sum_{j=1}^{n-1}v_{\rho,j}z^j\right)^l
                  \quad(2\le n<m).
 \end{split}
 \tag{CG11}\label{cg:inverse-recursion}
\end{equation}
To prove this, in the coefficient of $z^n$ in
$k_\rho(q_\rho(z))$, the unknown $v_{\rho,n}$ occurs only in the
linear term, with coefficient $d_\rho\ne0$. CG11 sets that
coefficient to zero for $n\ge2$, and the first line sets it to
one for $n=1$. It also follows that
\begin{equation}
 q_\rho(k_\rho(t))\equiv t\pmod{t^m}.
 \tag{CG12}\label{cg:two-sided}
\end{equation}
For a direct verification, substitution $z\mapsto k_\rho(t)$
is an invertible linear map on truncated polynomials: its matrix
is triangular with diagonal $1,d_\rho,\ldots,d_\rho^{m-1}$.
CG10 says that substitution by $q_\rho$ is a one-sided inverse;
finite-dimensional invertibility gives the other side. For $m=1$
the same assertion is the identity on constants, without a derivative
condition. Thus CG12
recovers $t$, not a replacement coordinate.

Set
\begin{equation}
 W_\rho(X)=\prod_{\sigma\ne\rho}(X-z_\sigma)^{m_\sigma},
 \qquad
 C_\rho(z)=\frac1{W_\rho(z_\rho+z)}\pmod{z^{m_\rho}}.
 \tag{CG13}\label{cg:reciprocal}
\end{equation}
The constant denominator is the exact nonzero product
$W_\rho(z_\rho)=\prod_{\sigma\ne\rho}(z_\rho-z_\sigma)^{m_\sigma}$.
For an explicit finite inverse, if
$W_\rho(z_\rho+z)=\sum_{j\ge0}w_{\rho,j}z^j$ and
$C_\rho(z)=\sum_{n=0}^{m_\rho-1}c_{\rho,n}z^n$, then
\begin{equation}
 c_{\rho,0}=w_{\rho,0}^{-1},\qquad
 c_{\rho,n}=-w_{\rho,0}^{-1}
              \sum_{j=1}^{n}w_{\rho,j}c_{\rho,n-j}
 \quad(1\le n<m_\rho).
 \tag{CG14}\label{cg:reciprocal-recursion}
\end{equation}
Coefficients of $W_\rho$ beyond its degree are zero. These are
exact finite polynomial operations at the given nodes.

\begin{theorem}[Complete inverse interpolation]\label{cg:inverse-thm}
For $f=(f_\rho(t))\in A_Z$, with
$f_\rho(t)=\sum_{j=0}^{m_\rho-1}f_{\rho,j}t^j$, define
\begin{equation}
 P_f(X)=\sum_{\rho\in Z}W_\rho(X)
    \left.
      \left[f_\rho(q_\rho(z))C_\rho(z)\right]_{<m_\rho}
    \right|_{z=X-z_\rho}.
 \tag{CG15}\label{cg:inverse}
\end{equation}
Then $\deg P_f<d$ and $P_f(H_c)=f$ in every full original local
factor. It is the unique polynomial of degree less than $d$ with
this property. In particular the exact inverse matrix is
\begin{equation}
 (E^{-1})_{k,(\rho,j)}=[X^k]\left\{
 W_\rho(X)
 \left.\left[q_\rho(z)^jC_\rho(z)\right]_{<m_\rho}
                                  \right|_{z=X-z_\rho}\right\},
 \quad 0\le k<d, 0\le j<m_\rho.
 \tag{CG16}\label{cg:inverse-matrix}
\end{equation}
For $j=0$, the factor $q_\rho(z)^0$ is one, also when $m_\rho=1$.
Moreover
\begin{equation}
 Q_H(X)=\prod_{\rho\in Z}(X-z_\rho)^{m_\rho}\in\mathbb{R}[X],\qquad
 \mathbb{R}[X]/(Q_H)\xrightarrow[\ X\mapsto H_c\ ]{\ \cong\ }A_{Z,\mathbb{R}},
 \quad
 \mathbb{C}[X]/(Q_H)\cong A_Z.
 \tag{CG17}\label{cg:real-generator}
\end{equation}
For $f\in A_{Z,\mathbb{R}}$ the polynomial $P_f$ in CG15 has real
coefficients.
\end{theorem}

\begin{proof}
Each summand in CG15 has degree at most
$(d-m_\rho)+(m_\rho-1)=d-1$.
At the node $z_\rho$, all other summands vanish modulo
$(X-z_\rho)^{m_\rho}$, while the $\rho$ summand has expansion
$f_\rho(q_\rho(z))$ modulo $z^{m_\rho}$ by CG13.
Substituting $z=k_\rho(t)$ and using CG12 gives exactly
$f_\rho(t)$ modulo $t^{m_\rho}$. This proves all original jets,
including the positive orders.
If a polynomial $P$ vanishes after substituting $H_c$ at $\rho$,
the invertible local substitution proved above says that
$P(z_\rho+z)$ vanishes modulo $z^{m_\rho}$.
Thus $(X-z_\rho)^{m_\rho}$ divides $P$.
For distinct nodes these powers are relatively prime; repeated
division, or the elementary factor theorem with multiplicities,
shows that their product $Q_H$ divides $P$.
A polynomial of degree less than $d$ divisible by $Q_H$ is zero.
This proves uniqueness and the complex isomorphism in CG17.
Applying CG15 to one coordinate vector $t_\rho^j$ gives CG16.
This is the specialized Hermite interpolation construction, with
the Chinese-remainder mechanism explicitly carried out; see the human-source
\href{https://stacks.math.columbia.edu/tag/00DT}{Stacks Project,
Lemma 10.15.4, tag 00DT} for the general theorem.

Since $H_c^\dagger=H_c$, one has
$z_{\iota\rho}=\overline{z_\rho}$, with the same multiplicity.
The coefficients of $Q_H$ are therefore real. If $P_f(H_c)=f$
and $f=f^\dagger$, then
$\overline{P_f}(H_c)=f^\dagger=f$.
Both polynomials have degree less than $d$, so uniqueness gives
$\overline{P_f}=P_f$. This proves surjectivity of the real map;
the same divisibility argument proves its kernel is $(Q_H)$.
It follows as well that $1,H_c,\ldots,H_c^{d-1}$ form a real
basis of $A_{Z,\mathbb{R}}$, and the same sequence is a complex
basis of $A_Z$.
\end{proof}

The auxiliary polynomial $Q_H$ is the newly derived relation in
the variable $X$. It does not replace $g(s)$, the original packet
polynomial $h(s)=c_h\prod(s-\rho)^{m_\rho}$, its chosen leading
coefficient $c_h$, or any original local unit $g/h$.
All of those data remain unchanged when CG15 is used as the
input vector of the original primitive lift.

For further direct checks on the real descent, let
$P_{\rho,j}$ denote the polynomial in braces in CG16.
Uniqueness and CG1 give
\begin{equation}
 P_{\iota\rho,j}(X)=(-1)^j\overline{P_{\rho,j}(X)}.
 \tag{CG18}\label{cg:paired-inverse}
\end{equation}
Indeed dagger sends the original jet vector $t_\rho^j$ supported
at $\rho$ to $(-1)^j t_{\iota\rho}^j$ supported at $\iota\rho$.
This proves CG18 and verifies every reflection sign in the
inverse formula.

\subsection{Finite Laurent words and every original prime factor}

In $\mathbb{C}[X_2^{\pm1},X_3^{\pm1}]$ define
\begin{equation}
 \mathcal H_c=A_+X_2+A_-X_2^{-1}+B_+X_3+B_-X_3^{-1},
 \quad A_\pm=\frac{1\mp ic}{2},\qquad
 B_\pm=\frac{c^2\mp ic^3}{2}.
 \tag{CG19}\label{cg:laurent-h}
\end{equation}
Under the original two-prime map $X_p\mapsto\lambda_p$, this is
exactly $H_c$. For $k\ge0$ its complete expansion is
\begin{equation}
 \mathcal H_c^k=
 \sum_{\substack{u_+,u_-,v_+,v_-\ge0\\u_++u_-+v_++v_-=k}}
 \frac{k! A_+^{u_+}A_-^{u_-}B_+^{v_+}B_-^{v_-}}
      {u_+!u_-!v_+!v_-!}
 X_2^{u_+-u_-}X_3^{v_+-v_-}.
 \tag{CG20}\label{cg:multinomial}
\end{equation}
This follows by choosing one of the four displayed terms in
each of the $k$ factors; the multinomial coefficient counts
exactly the choices with those four counts.

\begin{proposition}[A finite bound retaining all jets and original scales]
\label{cg:words-thm}
Every $f\in A_Z$ has the Laurent representative
$P_f(\mathcal H_c)$ supported on
\begin{equation}
 \{(u,v)\in\mathbb{Z}^2:|u|+|v|\le d-1\}.
 \tag{CG21}\label{cg:diamond}
\end{equation}
There are exactly $2d^2-2d+1$ possible distinct Laurent words
in this set. If $f\in A_{Z,\mathbb{R}}$ then $P_f$ is real and the
Laurent coefficient at $(-u,-v)$ is the complex conjugate of
the coefficient at $(u,v)$.
For every integer pair, its action on the original packet is
\begin{equation}
 X_2^uX_3^v\longmapsto
 M_{2^{u(s-1/2)}3^{v(s-1/2)}}
 =2^{-u/2}3^{-v/2}\Pi_2^u\Pi_3^v,
 \qquad u,v\in\mathbb{Z}.
 \tag{CG22}\label{cg:original-words}
\end{equation}
For any further rational prime $p$, define the entire functions
$a_p,b_p$ by CG2. Their full jets have unique real polynomials
$P_{a_p},P_{b_p}$ given by CG15, each of degree at most $d-1$.
Then its original action is exactly
\begin{equation}
 \Pi_p=p^{1/2}M_{\,P_{a_p}(H_c)+iP_{b_p}(H_c)}.
 \tag{CG23}\label{cg:all-primes}
\end{equation}
In particular the same finite Laurent support bound suffices
for every other original prime action on the given full packet.
\end{proposition}

\begin{proof}
Each monomial in CG20 has
$|u_+-u_-|+|v_+-v_-|\le k$.
Since $P_f$ has degree at most $d-1$, its expansion is supported
in CG21. For $D\ge0$, the shell $|u|+|v|=j$ has $4j$ points
for $j\ge1$, while the central shell has one. The total is
$1+\sum_{j=1}^D4j=1+2D(D+1)$; setting $D=d-1$ gives the
claimed exact cardinality. The Laurent involution conjugates
coefficients and sends $X_p$ to $X_p^{-1}$.
In CG19, $A_- =\overline{A_+}$ and
$B_-=\overline{B_+}$, so $\mathcal H_c$ is fixed.
A real polynomial in it is also fixed, proving the coefficient
relation.

The original operators $\Pi_p$ are invertible on $A_Z$ because
the constant term $p^\rho$ is nonzero at every root.
Thus $X_p=p^{-1/2}\Pi_p$ and
$X_p^{-1}=p^{1/2}\Pi_p^{-1}$ hold with their full jets.
Taking integer powers proves CG22, including negative exponents.
Finally $a_p+ib_p=\lambda_p$ as an entire identity, so CG15
applied to its two real components recovers the full Taylor
series of $\lambda_p$ at every root. Multiplication by the
retained original factor $p^{1/2}$ proves CG23.
\end{proof}

One may equivalently apply CG15 directly to
\begin{equation}
 (\lambda_p)_\rho(t)=p^{\rho-1/2}
              \sum_{j=0}^{m_\rho-1}\frac{(\log p)^j}{j!}t^j,
 \qquad
 (p^s)_\rho(t)=p^{\rho}
              \sum_{j=0}^{m_\rho-1}\frac{(\log p)^j}{j!}t^j.
 \tag{CG24}\label{cg:other-prime-jets}
\end{equation}
Linearity of CG15 then gives
$P_{p^s}=p^{1/2}P_{\lambda_p}$ and
$P_{\lambda_p}=P_{a_p}+iP_{b_p}$.
No prime eigenvalue has been replaced by its modulus.

\subsection{Edge cases and the exact scope of the bounds}

The case $r=1,m_\rho=1$ gives $d=1$, $E=(1)$,
$W_\rho=1$, $q_\rho=0$, and $P_f=f_\rho$.
The empty determinant products in CG8 equal one.
CG21 consists only of $(0,0)$ and has cardinality one.
The selector CG4 gives $B=0$ and permits $c=0$ without a
derivative test.
Since a one-root reflection-stable packet is fixed by reflection,
its real algebra in this case is $\mathbb{R}$.

For $r=1$ with any larger multiplicity, CG4 gives $B=3$.
Distinct-node conditions
are empty and a nonzero derivative still gives the full truncated
factor. The determinant is exactly
$d_\rho^{m_\rho(m_\rho-1)/2}$. The local recursion retains each
of its $m_\rho-1$ positive-order coefficients.
No argument assumes the multiplicities are equal at different
orbits; equality is needed only at the two members of one
reflection orbit, as provided by the original reflection.

A chosen integer may be zero. In that case $B_\pm=0$, and terms
with positive corresponding multinomial counts vanish. This can
only reduce Laurent support; it does not invalidate the bound.
Some $z_\rho$ may be zero. The determinant and inverse use only
differences of distinct nodes and nonzero $d_\rho$ at nonsimple
roots, so this
causes no division by zero. Interpret $z^0=1$ and constant
polynomial coefficients directly, including the $m_\rho=1$
case in CG16.

The degree bound $d-1$ is optimal for representing every element
by polynomials in this one generator: polynomials of degree at
most $D<d-1$ have dimension $D+1<d$, so cannot span $A_Z$.
No optimality assertion is made for arbitrary two-variable
Laurent representatives. The $2d^2-2d+1$ words are a spanning
list with many relations, not a basis of a $d$-dimensional
space. Products of two unreduced representatives may have
Laurent degree $2d-2$, but their packet product has a new
representative of degree at most $d-1$ by CG15 or reduction
modulo $Q_H$; no original jet is lost in this reduction.

The candidate count in CG4 depends on the number of distinct
roots, whereas the interpolation and Laurent-degree counts
depend on the sum of full multiplicities. Replacing $d$ by $r$
in those latter bounds would discard the nilpotent directions
when a multiplicity exceeds one.
The finite selector produces a generator for each supplied
finite packet. It does not assert the existence of a single
integer that works simultaneously for all finite packets, nor
does it assert a bound independent of the packet size.

